\section{Methodology}

\subsection{Actionable Specificity Framework}

We propose that feedback effectiveness for LLM code repair depends on \textbf{Actionable Specificity (AS)}---the degree to which a verification signal provides information that enables targeted bug localization and fix inference. We decompose AS into three orthogonal dimensions:

\textbf{AS\_loc (Localization):} Whether the signal identifies the failure location. Binary: 1 if file:line information present, 0 otherwise. Extraction pattern: \texttt{File "([{\textasciicircum}"]+)", line (\textbackslash d+)}.

\textbf{AS\_state (State Exposure):} How many variable values are visible at the failure point. Continuous: count of variable-value pairs. Extraction pattern: \texttt{(\textbackslash w+)\textbackslash s*=\textbackslash s*(["']?[\textbackslash w\textbackslash d\textbackslash.\textbackslash-\textbackslash[\textbackslash]\{\}]+["']?)}.

\textbf{AS\_causal (Causal Context):} Execution trace depth between variable assignment and failure. Continuous: count of traceback frames. Extraction pattern: \texttt{{\textasciicircum}\textbackslash s+File "([{\textasciicircum}"]+)", line (\textbackslash d+), in (\textbackslash w+)}.

These components are designed to be \textbf{independently measurable} from signal text using regex extraction, enabling quantitative analysis across signal types.

\subsection{Signal Generation Conditions}

We defined 6 signal generation conditions varying in expected AS level:

\begin{table}[h]
\centering
\begin{tabular}{@{}lll@{}}
\toprule
Condition & Description & Expected AS Level \\
\midrule
C1 & Full runtime trace & High (all components) \\
C2 & Truncated trace (last 5 frames) & Medium-High \\
C3 & Value-masked trace & Medium (no AS\_state) \\
C4 & Error message only & Low \\
C5 & Static analysis (mypy/pylint) & Variable \\
C6 & Syntax error baseline & Minimal \\
\bottomrule
\end{tabular}
\end{table}

The conditions are ordered such that C1 $\geq$ C2 $\geq$ C3 $\geq$ C4 in expected AS, enabling ordered hypothesis testing.

\subsection{Design Rationale}

We chose the \textbf{simplest operationalization} (regex extraction) to establish feasibility bounds before investing in more sophisticated methods (AST analysis, LLM-based extraction). This follows the scientific principle of testing necessary conditions first: if AS components cannot be reliably extracted via simple patterns, the framework requires reformulation before testing predictive power.

\subsection{Existence Test Design (H-E1)}

Before testing whether AS predicts repair success, we test whether AS is \emph{measurable}. Our existence hypothesis:

\begin{quote}
\textbf{H-E1:} Under standard verification signal generation, AS\_loc, AS\_state, and AS\_causal can be independently computed for $\geq$95\% of signals, and AS values show expected ordering across conditions (C1$\geq$C2$\geq$C3$\geq$C4).
\end{quote}

\textbf{Success Criteria:}
\begin{itemize}
    \item Primary: Extraction rate $\geq$95\% across all signals
    \item Secondary: AS ordering C1$\geq$C2$\geq$C3$\geq$C4 preserved
\end{itemize}
